% ZHOU-SOURCE-PAGE: 5 PRINTED: FRONTMATTER
\begin{center}
\emph{To the memory of my sister, Xinli Zhou}
\end{center}

\textcopyright{}Copyright 2008 by Xinfeng Zhou, \url{http://www.quantfinanceinterviews.com}

All right reserved.

No part of this book may be reproduced or transmitted in any form or by any means,
electronic or mechanical, including photocopying, recording or by any information
storage and retrieval system, without the written permission of the Publisher, except
where permitted by law.

\clearpage
% ZHOU-SOURCE-PAGE: 13 PRINTED: FRONTMATTER
\section*{Preface}
\addcontentsline{toc}{section}{Preface}

This book will prepare you for quantitative finance interviews by helping you zero in on
the key concepts that are frequently tested in such interviews. In this book we analyze
solutions to more than 200 \emph{real} interview problems and provide valuable insights into
how to ace quantitative interviews. The book covers a variety of topics that you are
likely to encounter in quantitative interviews: brain teasers, calculus, linear algebra,
probability, stochastic processes and stochastic calculus, finance and programming.

Professionals and students seeking to pursue a career in quantitative finance or related
quantitative fields will benefit most from thoroughly reading this book. In recent years,
we have seen a dramatic surge in demand for talents with strong quantitative skills from
investment banks, investment management firms, hedge funds, financial software
vendors and financial consulting companies. As a result, quant, an umbrella description
that encompasses quantitative analysts, quantitative researchers, quantitative strategists,
quantitative traders, and quantitative developers, has become an attractive career choice.

Dozens of financial engineering or computational finance programs have been
established in the last few years to educate professionals for quantitative finance jobs.
Graduates with backgrounds in finance, mathematics, physics, computer sciences, and
various engineering majors are contending for quant jobs as well. Naturally, the
competition is fierce. To be a successful candidate, you have to distinguish yourself
from many other excellent applicants.

In general, a successful candidate for a quantitative finance position is expected to have
a strong mathematics background (in probability, statistics, stochastic calculus, etc.),
solid programming skills and basic to intermediate-level finance knowledge. Most
candidates find quantitative interviews, or at least some interview problems, challenging.

Quantitative interviews cover a broad range of mathematics, finance and programming
topics that the candidates may have never used or even encountered in their daily work
or study. Moreover, most interview problems require strong problem-solving skills,
beyond reciting formulas or doing simple calculations. A successful candidate needs a
combination of knowledge and problem-solving skills in order to excel in quantitative
interviews. This is precisely what this book provides!

This book addresses these aspects by reviewing the necessary finance and mathematical
concepts that serve as tools to structure and solve interview problems. Since it includes
most of the topics used by quantitative interviewers, it presupposes some basic
preparation in mathematics, statistics, finance, and programming.

I also strongly recommend that you try to solve each problem on your own first before
reading the answer. Working out solutions on your own will help you improve your
problem-solving skills and help you quickly identify common approaches to tackling
quantitative problems.

% ZHOU-SOURCE-PAGE: 14 PRINTED: FRONTMATTER
Needless to say, you are likely to encounter some problems in interviews that are similar
to or exactly the same as the problems in this book. After all, the book covers many
essential quantitative topics using real interview problems. However, the goal of the
book is not to teach you how to game the system by remembering the answers! In fact,
just memorizing answers may not help much in your interview process. Unless you truly
understand the underlying concepts and can analyze the problems yourself, you will fail
to elaborate on the solutions and will be ill-equipped to answer many other problems
that use similar concepts. (Besides, many experienced quantitative interviewers are good
at catching those who have simply memorized ``canned'' answers.)

This is exactly the reason why I make significant effort to review essential concepts, to
present solution strategies, and to analyze the solutions in detail instead of simply
providing answers to problems. Furthermore, although the building blocks can be
learned, how one analyzes problems and implements these concepts usually makes a big
difference---and these are the skills you can acquire through practice, practice and
practice.

I realize that there may be better methods to solve some of the problems presented in
this book. It is entirely possible that despite my best efforts some inadvertent errors may
have crept in. Please email me at xinfeng@quantfinanceinterviews.com if you have a
better approach to solving some of these problems or find errors. I will be grateful for
your feedback and will post corrections and your constructive feedback on the book's
companion website \url{http://www.quantfinanceinterviews.com}. The website is a joint
venture with my editor, Brett Jiu. You will also find some extra interview problems
with answers that we have gathered.

I sincerely hope that you enjoy solving these problems and are successful in your
interviews.

\begin{flushright}
Xinfeng Zhou
\end{flushright}

% ZHOU-SOURCE-PAGE: 15 PRINTED: FRONTMATTER
\section*{Notations}
\addcontentsline{toc}{section}{Notations}

\begin{tabular}{@{}ll@{}}
  $\forall$ & for each/for every/for all \\
  $\exists$ & there exists \\
  $\therefore$ & therefore \\
  $A \Rightarrow B$ & whenever $A$ is true, $B$ is also true \\
  \textit{s.t.} & such that \\
  $a \wedge b$ & the minimum of $a$ and $b$ \\
  $a \vee b$ & the maximum of $a$ and $b$ \\
  $\displaystyle\sum_{i=1}^{n} x_i$ & $x_1+x_2+\cdots+x_n$ \\
  $\displaystyle\prod_{i=1}^{n} x_i$ & $x_1\cdot x_2\cdot\cdots\cdot x_n$ \\
  $n!$ & factorial of nonnegative integer $n$, $n!=\displaystyle\prod_{i=1}^{n}i$ ($0!=1$) \\
  $x\%y$ & modulo operation \\
  $\Phi$ & empty set \\
  $\displaystyle\int f(x)\,\dd x$ & indefinite integral of $f(x)$ \\
  $\displaystyle\int_a^b f(x)\,\dd x$ & definite integral of $f(x)$ from $a$ to $b$ \\
  $x^+$ & $\max(x,0)$ \\
  $N(\mu,\sigma^2)$ & normal distribution with mean $\mu$ and variance $\sigma^2$ \\
  $\CDF$ & cumulative density function \\
  $\PDF$ & probability density function
\end{tabular}
